\begin{table*}[!t]
\centering
\caption{Module~10 constants.}
\label{tab:m10const}
\footnotesize
\renewcommand{\arraystretch}{1.15}
\begin{tabularx}{\textwidth}{@{}l r l c >{\raggedright\arraybackslash}X@{}}
\toprule
Symbol & Value & Unit & Basis & Meaning, justification, source\\
\midrule
$[\mathrm{RegA}]_\mathrm{tot}$ & \num{1540} & $\mathrm{molec/voxel}$ & E & Copy number never measured; $0.32\,\mu$M puts RegA on the scale of ERK2 and PKA and keeps the required catalytic efficiency sub-diffusion-limited. \\
$[\mathrm{PKA}]_\mathrm{tot}$ & \num{1540} & $\mathrm{molec/voxel}$ & E & $0.32\,\mu$M. \\
$k_{B1}$ & \num{0.05} & $\mathrm{s^{-1}}$ & E & \mkt{M10-relay}{Lumped phospho-relay (RdeA$\to$RegA D212) that activates RegA}~\cite{thomason1999,thomason1998}. \\
$k_{B2}$ & \num{5.6e-3} & $\mathrm{s^{-1}}$ & C & \mkt{M10-relay}{Reverse phospho-transfer}~\cite{thomason1999}; chosen for a resting active fraction of 0.9. \\
$k_{B3}$ & \num{8} & $\mu\mathrm{M}^{-1}\mathrm{s}^{-1}$ & C & \mkt{M10-erk-rega}{ERK2$^*$ phosphorylates and thereby inhibits RegA}~\cite{laub1998,maeda2004}. The Laub--Loomis coefficient (\mkt{M10-laub-k8}{$k_8=1.3$ in relative units}~\cite{laub1998}) is not a molecular rate; the value gives $\tau\approx0.5$\,s at the ERK2$^*$ peak so that the brake is released inside the X spike. \\
$k_{B4}$ & \num{0.0167} & $\mathrm{s^{-1}}$ & C & The reset reaction ($\tau=60$\,s) that ends each pulse; of the order of the RegA replenishment coefficient \mkt{M10-laub-k7}{$k_7=2\,\mathrm{min^{-1}}$ of the Laub--Loomis model}~\cite{laub1998} and set so that the cell has reset within about 6\,min. \\
$k_\mathrm{cat}$ & \num{30} & $\mu\mathrm{M}^{-1}\mathrm{s}^{-1}$ & D & $k_\mathrm{cat}/K_M=3\times10^{7}\,\mathrm{M^{-1}s^{-1}}$: $k_\mathrm{cat}\approx150\,\mathrm{s^{-1}}$ at \mkt{M10-rega-km}{the measured $K_M\approx5\,\mu$M}~\cite{thomason1998}. \\
$\phi$ & \num{0.05} &  & L & Activity of RegA$_u$ relative to RegA$_p$; \mkt{M10-rega-20}{phosphorylation activates RegA by at least 20-fold}~\cite{thomason1999}. \\
$k_{D1}$ & \num{0.25} & $\mu\mathrm{M}^{-1}\mathrm{s}^{-1}$ & E & cAMP binding to the PKA regulatory subunit; with $k_{D2}$ this gives $K_d=k_{D2}/k_{D1}=0.1\,\mu$M (model choice; \mkt{M10-circuit}{PKA activation is the second element of the circuit of}~\cite{laub1998}). \\
$k_{D2}$ & \num{0.025} & $\mathrm{s^{-1}}$ & D & PKA$^*$ decay ($\tau=40$\,s), from \mkt{M10-laub-k4}{the Laub--Loomis coefficient $k_4=1.5\,\mathrm{min^{-1}}$}~\cite{laub1998}; sets the minutes-long refractory tail. \\
$k_{D3}$ & \num{0.0417} & $\mu\mathrm{M}^{-1}\mathrm{s}^{-1}$ & D & PKA$^*$ inhibits ERK2, from \mkt{M10-laub-k6}{the coefficient $k_6=0.8$ of}~\cite{laub1998} (relative units, normalised here by the PKA pool); PKA also \mkt{M10-pka-erk}{shapes ERK2 activation and adaptation}~\cite{aubry1997}. \\
$D_{\mathrm{RegA}},D_{\mathrm{PKA}},D_{\mathrm{ERK2}}$ & \num{10} & $\mu\mathrm{m^2/s}$ & E & Cytosolic value. \\
\bottomrule
\end{tabularx}
\end{table*}
